\chapter{Vision Loss}

\section*{Definition}
Vision loss is partial or complete reduction of sight in one or both eyes. It may be sudden or gradual and may originate in the eye, optic nerve, brain, or visual blood supply.

\section*{What Happens in Your Body}
The symptom results from irritation, inflammation, altered organ function, injury, infection, or disruption of normal neurologic or vascular signaling. Age, pregnancy status, onset, distribution, severity, and associated findings determine the urgency and likely causes.

\section*{Causes}
\begin{itemize}
  \item Refractive or lens disease, glaucoma, retinal disease, optic neuritis, vascular occlusion, stroke, trauma, infection, and inflammatory disease.
  \item Medication effects, environmental exposures, and underlying chronic disease may also contribute.
  \item Serious causes are less common but must be considered when symptoms are sudden, progressive, severe, or associated with systemic illness.
\end{itemize}

\section*{When to See a Doctor}
Sudden vision loss, severe eye pain, a curtain or shadow, flashes with floaters, or vision loss with weakness, speech difficulty, confusion, or severe headache is an emergency. Arrange prompt evaluation for persistent, recurrent, unexplained, or function-limiting symptoms.

\section*{Related Symptoms}
\begin{itemize}
  \item Pain, fever, fatigue, nausea, weakness, or changes in normal function
  \item Bleeding, swelling, discharge, breathing difficulty, or neurologic symptoms when a serious cause is present
\end{itemize}

\section*{Self-Care / Home Management}
Avoid known triggers, protect the affected area, maintain appropriate hydration, and record timing and associated symptoms. Use only age- or pregnancy-appropriate medicines recommended by a clinician. Do not delay emergency care while trying home treatment.

\section*{How Your Doctor Will Evaluate You}
\begin{itemize}
  \item History addresses onset, duration, severity, triggers, injuries, exposures, medications, medical conditions, age, and pregnancy status when relevant.
  \item Examination includes vital signs and focused assessment of the relevant organ system, neurologic status, hydration, and function.
  \item Testing may include blood or urine studies, cultures, imaging, physiologic testing, fetal or pediatric assessment, or specialist examination.
\end{itemize}

\section*{Treatment Options}
Treatment is directed at the cause and may include supportive care, targeted medication, treatment of infection or inflammation, correction of metabolic disease, physical therapy, or specialist procedures. Persistent symptoms require reassessment.

\section*{References}
\begin{thebibliography}{99}
\bibitem{symptom_vision_loss:ref1} National Institute for Health and Care Excellence. Relevant clinical guidance. NICE; accessed 2025.
\bibitem{symptom_vision_loss:ref2} American Academy of Pediatrics or relevant specialty society guidance. Current clinical recommendations; accessed 2025.
\bibitem{symptom_vision_loss:ref3} NHS. Symptoms A to Z and related clinical information. Accessed 2026.
\end{thebibliography}
